% Options for packages loaded elsewhere
\PassOptionsToPackage{unicode}{hyperref}
\PassOptionsToPackage{hyphens}{url}
%
\documentclass[
  11pt,
]{article}
\usepackage{amsmath,amssymb}
\usepackage{iftex}
\ifPDFTeX
  \usepackage[T1]{fontenc}
  \usepackage[utf8]{inputenc}
  \usepackage{textcomp} % provide euro and other symbols
\else % if luatex or xetex
  \usepackage{unicode-math} % this also loads fontspec
  \defaultfontfeatures{Scale=MatchLowercase}
  \defaultfontfeatures[\rmfamily]{Ligatures=TeX,Scale=1}
\fi
\usepackage{lmodern}
\ifPDFTeX\else
  % xetex/luatex font selection
\fi
% Use upquote if available, for straight quotes in verbatim environments
\IfFileExists{upquote.sty}{\usepackage{upquote}}{}
\IfFileExists{microtype.sty}{% use microtype if available
  \usepackage[]{microtype}
  \UseMicrotypeSet[protrusion]{basicmath} % disable protrusion for tt fonts
}{}
\makeatletter
\@ifundefined{KOMAClassName}{% if non-KOMA class
  \IfFileExists{parskip.sty}{%
    \usepackage{parskip}
  }{% else
    \setlength{\parindent}{0pt}
    \setlength{\parskip}{6pt plus 2pt minus 1pt}}
}{% if KOMA class
  \KOMAoptions{parskip=half}}
\makeatother
\usepackage{xcolor}
\usepackage[margin=1in]{geometry}
\setlength{\emergencystretch}{3em} % prevent overfull lines
\providecommand{\tightlist}{%
  \setlength{\itemsep}{0pt}\setlength{\parskip}{0pt}}
\setcounter{secnumdepth}{-\maxdimen} % remove section numbering
\ifLuaTeX
  \usepackage{selnolig}  % disable illegal ligatures
\fi
\usepackage{bookmark}
\IfFileExists{xurl.sty}{\usepackage{xurl}}{} % add URL line breaks if available
\urlstyle{same}
\hypersetup{
  pdftitle={Ramified quadratic transport of ordinary zero-Hessian currents},
  pdfauthor={Xean research note},
  hidelinks,
  pdfcreator={LaTeX via pandoc}}

\title{Ramified quadratic transport of ordinary zero-Hessian currents}
\author{Xean research note}
\date{3 October 2026}

\begin{document}
\maketitle

\subsection{1. Exact statement and
status}\label{exact-statement-and-status}

Let \(K\) be a complete nontrivially valued non-Archimedean field and
let \(E/K\) be a separable quadratic extension. Write \(\sigma\) for its
nonidentity automorphism and \(\pi:X_E\to X\) for scalar projection. No
condition on the ramification, residue fields, characteristic, or size
of \(a-\sigma(a)\) is imposed.

\textbf{Theorem.} Suppose that \(X\) is good, separated, strict,
boundaryless, and pure of analytic dimension \(n\geq1\), with its actual
analytic structure, possibly singular and nonreduced. For every
continuous finite \(h:X\to\mathbb R\), \[
 d'd''[h]=0\quad\Longleftrightarrow\quad
 d'd''[h\circ\pi]=0.                                      \tag{1}
\] Both equations use ordinary CLD compactly supported \((n-1,n-1)\)
tests and the canonical calibrations for their respective fields. This
applies to every ordinary open subspace and iterates along finite towers
of separable quadratic extensions. In particular it proves the forward
implication for \(K=\mathbb Q_2\), \(E=\mathbb Q_2(\sqrt2)\) on every
space in this scope, including every pure singular or nonreduced
analytic curve. No upstairs stability of the equation is assumed.

The new assertion is the forward implication without a
residue-separating unit. The finite-separable reverse implication alone
is not presented as new progress. The theorem supplies a scalar-current
transport gate; it does not supply classical recognition or ordinary
two-sign psh regularization. Thus the original atom P06.2 remains open
and the status is \texttt{partial}, with \texttt{candidate=false}. This
partial theorem received an independent source-aware review through
Xean native Codex, using gpt-6.1-sol with xhigh reasoning. The complete
proof and review are preserved with exact artifact bindings. This
review is not formal verification.

For a unit \(f\), set \(t_f=-\log|f|\). We use the ordinary form sheaf
\(A^{p,q}\), not the weakly smooth GJR sheaf. In a pure
\(n\)-dimensional space the zero equation can be tested as \[
 B_h(\alpha)=\int_X h\,(d'd''\alpha)_{\mathrm{Rad},K}=0
 \quad(\alpha\in A_c^{n-1,n-1}(X)).                         \tag{2}
\] The fixed convention-dependent duality sign does not change
vanishing. All compact sets, neighbourhoods and supports refer to the
ordinary Hausdorff topology. An affinoid neighbourhood contains an
ordinary open neighbourhood; its topological interior in \(X\) is
distinct from its absolute analytic interior.

The proof does \textbf{not} assert exact descent of invariant polynomial
superforms in the ramified case. It constructs ordinary downstairs tests
whose pulled-back differentiated measures approximate every averaged
upstairs test in total variation, with fixed compact support. This
distinction is essential at the fold.

\subsection{2. Exact primary-source
inputs}\label{exact-primary-source-inputs}

The numbering here refers to the cached primary snapshots
Chambert-Loir--Ducros (CLD), arXiv:1204.6277v2, 25 July 2025, and
Gubler--Rabinoff (GR), arXiv:2506.13548v2, 17 September 2025. The
following are the precise source inputs. They are stated here so that
the proof can be read without opening a local source file. The
elementary new algebraic, analytic and approximation lemmas are proved
below. No earlier manuscript or catalogue recognition result is a
premise.

\begin{enumerate}
\def\labelenumi{\arabic{enumi}.}
\item
  \textbf{Affinoid foundations.} CLD 11.1.1--2: affinoid spectra are
  compact Hausdorff; affinoid algebras are noetherian, ideals are
  closed, and \(\mathcal M(A)\to\operatorname{Spec}A\) is surjective.
  CLD 11.1.5 supplies rational affinoid neighbourhoods inside prescribed
  ordinary opens in an affinoid. CLD 11.1.7: good spaces have affinoid
  neighbourhood bases; ordinary opens are analytic domains. CLD 11.7.7:
  over a nontrivially valued field, every point of a good strict space
  has a strict affinoid neighbourhood. CLD 11.1.9: affinoid fibre
  products are computed by completed tensor products; scalar projection
  is surjective. CLD 11.4.4: a finite algebra over an affinoid algebra,
  with its finite-module Banach topology, is affinoid. We prove the
  particular finite tensor and shrinking assertions below.
\item
  \textbf{Dimension and ambient structure.} CLD 11.2.2 says analytic
  separatedness implies topological separatedness. CLD 11.2.3 says
  strict affinoid dimension is Krull dimension, scalar extension
  preserves dimension, analytic domains preserve local analytic
  dimension, and for the minimal-prime components \(X_i\) of an
  affinoid, \(\dim_xX=\sup_{x\in X_i}\dim X_i\). GR 2.11(5') states, for
  any analytic extension \(L/K\) and analytic morphism \(Y\to Z\),
  \(\partial(Y_L/Z_L)\subseteq\pi^{-1}(\partial(Y/Z))\). In particular a
  boundaryless space remains boundaryless. We check minimal components,
  strictness, goodness and separatedness below; equality of total
  dimensions is not substituted for purity.
\item
  \textbf{Ordinary moment descriptions and nullity.} CLD Proposition
  15.1.1(b): on a topologically separated analytic space every ordinary
  \((p,q)\) form has an ordinary open cover on each member of which it
  is a unit-moment pullback of a coefficient form on a paralinéaire set
  containing the moment image. CLD 3.3.3 says that such a coefficient
  form can be represented by a Euclidean form on an ambient open. It
  uses Lemma 1.9.4 in its precise closed-subset form: if a paralinéaire
  subset \(P\) is closed in an ambient open \(U\), a coefficient form on
  \(P\) extends to \(U\); 3.3.3 chooses such a \(U\). CLD 6.2.1--2
  defines ordinary smooth forms as G-form sections having these
  descriptions on ordinary neighbourhoods. They form an ordinary sheaf
  stable under wedge, \(J,d',d''\). Pullbacks commute with these
  operations (CLD 3.3.12). CLD Proposition 3.4.4 states: for a
  G-tropical chart \(f:X\to V\), a coefficient form \(\eta\) has zero
  pullback in the G-form sheaf if and only if its restriction to
  \(f(W)\) is zero for every compact domain \(W\) of \(X\). Proposition
  3.4.1 is the compact-chart presheaf version. We use actual compact
  images to prove the tripod one-form identity. Its moment and
  coefficients are ordinary, so its equality in the containing G-sheaf
  is an equality of ordinary forms.
\item
  \textbf{Cutoffs.} CLD 15.2.2: a good topologically separated analytic
  space is locally holomorphically separated. For a compact subset \(F\) and a finite ordinary open cover
  \((U_i)\) of \(F\), there are nonnegative ordinary smooth
  \(\theta_i\), each compactly supported inside \(U_i\), with
  \(\sum_i\theta_i=1\) near \(F\). This compact specialization follows
  from CLD Corollary 15.2.5 and ordinary smooth functional calculus.
  For each \(x\in F\), apply 15.2.5 to \(\{x\}\subset U_i\) to obtain a
  smooth bump equal to one near \(x\) and compactly supported in
  \(U_i\). Compose with a nonnegative smooth real function equal to
  zero near zero and one near one. Select finitely many such bumps by
  compactness and group them as \(\psi_i\). Set \(s=\sum_i\psi_i>0\)
  near \(F\). Apply 15.2.5 again to obtain a nonnegative smooth
  \(\chi=1\) near \(F\), compactly supported inside \(\{s>0\}\).
  On that open set put \(\theta_i=\chi\psi_i/s\), extended by zero.
  This extension is smooth because the support of \(\chi\) lies
  compactly inside the open set. Every \(\theta_i\) has compact support
  in \(U_i\), and their sum is \(\chi\). Only this compact, finite-cover
  statement is used here. Extension by zero of products with these cutoffs is proved
  directly below, without global paracompactness.
\item
  \textbf{Canonical calibrated top measures.} CLD Theorem 15.7.1: a
  topologically separated pure \(n\)-dimensional analytic space with its
  canonical calibration and analytic boundary is a tropical space of
  level \(n\). CLD 15.7.2: the induced structure on an analytic domain
  agrees with its own analytic structure. CLD 5.4.4: every compact
  tropical chart is pleasant and its tropical polyhedron is the
  top-dimensional part of the compact image. CLD 14.3.2--5 constructs
  the positive standard canonical calibration on skeletal cells. CLD
  4.5.1--2 says a standard calibration is strictly positive and
  G-locally constant; on a paralinéaire subset of an affine space it has
  a cell decomposition on each top cell of which it is a constant volume
  vector. CLD 14.5.1 defines the moment-image calibration as the direct
  image of the canonical skeletal calibration and says that it is
  standard. CLD 14.5.2 gives its explicit value on a top cell \(D\):
  choose a torus projection \(q\) injective on the cell, let \(\tau_D\)
  be the affine inverse onto its span, and let \(d_D\) be the sum of the
  finite-flat moment degrees \(d_C\) of the skeletal cells above \(D\).
  Then \(\mu_{\Pi_f,D}=d_D\,\tau_{D,*}|e_1\wedge\cdots\wedge e_n|\).
  This is the constant finite positive cell density used in Section 5.
  Its degrees are degrees of actual finite fibre algebras, including
  their lengths. CLD 5.8.3--5 defines a top form's Radon measure by
  pairing the form with this calibration on its supporting skeleton. CLD
  Proposition 5.10.1(c) states: for a pleasant chart \(f\) and a Borel
  \((p,n)\) coefficient form \(\eta\), \(p\in\{n-1,n\}\), the associated
  Borel measure is locally finite precisely when the restricted
  coefficient form is locally integrable on the calibrated tropical
  polyhedron; then \(f_*((f^*\eta)_{\mathrm{Rad}})
      =\langle\eta,\mu_{\Pi_f}\rangle_{\mathrm{Rad}}\). Part (d), for
  \(p=n\), gives
  \(\int_X f^*\eta=\int_{\Pi_f}\langle\eta,\mu_{\Pi_f}\rangle\). Thus a
  smooth top form is finite on a compact moment domain. Lower
  dimensional cells carry no top measure. CLD 7.3.5 defines integration
  currents for locally integrable multipliers; continuous finite
  multipliers are locally integrable. The matrix estimate and fold
  total-variation bound needed here are proved in Sections 5 and 7.
\item
  \textbf{Scalar calibration, without an extension-degree factor.} CLD
  Proposition 14.3.6: if \(\Sigma\) is a skeletal subset of a
  level-\(n\) analytic space, \(L/K\) complete, and \(\Sigma_L\) its
  base change with canonical calibration \(\mu_L\), then
  \(\pi_*\mu_L=\mu\). CLD Proposition 15.3.1 consequently states, for a
  \((n,n)\) G-form \(\omega\) with tropically Borel coefficients,
  \(\pi_*\omega_{L,\mathrm{Bor}}=\omega_{\mathrm{Bor}}\), and for a
  locally integrable such form, \[
    \pi_*\omega_{L,\mathrm{Rad}}=\omega_{\mathrm{Rad}}.       \tag{3}
  \] We use this for smooth differentiated forms. The proof of 14.3.6
  computes the degree of a convenient finite-flat moment projection as
  the sum of the degrees at its upstairs skeletal lifts (CLD Lemma
  13.3.3), retaining fibre-algebra dimensions. Its normalization is (3),
  not \([L:K]\) times (3). No pullback of a continuous current is an
  input.
\item
  \textbf{Same-carrier restriction of field and isomorphisms.} CLD
  Remark 14.4.3: if \(F/K\) is finite and \(Z\) is an \(F\)-analytic
  space of level \(n\), its restricted \(K\)-analytic structure \(Z_0\)
  has the same dimension and skeletal subsets, and their canonical
  calibrations satisfy \(\mu_{Z_0}=[F:K]\mu_Z\). We prove equality of
  the ordinary form sheaves on this same carrier in Section 3. CLD
  15.4.1 and Proposition 15.4.3 state: for a map of level-\(n\) analytic
  spaces finite away from a Zariski closed subset of dimension \(<n\),
  and a locally integrable \((p,n)\) G-form, \(p\in\{n-1,n\}\),
  \(f_*((f^*\omega)_{\mathrm{Rad}})
      =\mathrm{dc}(f)\,\omega_{\mathrm{Rad}}\). We use only a
  \(K\)-analytic isomorphism, of calibrated degree one, on \(Z_0\).
  After cancelling the restriction factor two this justifies semilinear
  averaging over \(E\), as distinct from scalar projection.
\end{enumerate}

These source statements apply to the singular and nonreduced spaces
specified here. None requires normalizing \(X\) or discarding generic
lengths. No weakly smooth descent theorem, analytic-morphism descent
theorem, psh recognition theorem, or manuscript receipt is used as an
ordinary form/current theorem.

\subsection{3. Finite tensors, fibres, the involution, and
purity}\label{finite-tensors-fibres-the-involution-and-purity}

\subsubsection{3.1. Actual algebra and
shrinking}\label{actual-algebra-and-shrinking}

If the minimal polynomial is \(Z^2+pZ+q\), its other root is
\(-p-a\in E\) and is distinct from \(a\) by separability. Thus every
separable quadratic extension is Galois with the stated involution.
Choose \(a\in E\setminus K\) and put \(c=a-\sigma(a)\ne0\). \(1,a\) is a
\(K\)-basis and \(\sigma(c)=-c\), including in characteristic two. For a
\(K\)-affinoid Banach algebra \(A\), \[
 B=A\widehat\otimes_K E=A\otimes_K E=A\oplus Aa.             \tag{4}
\] Here is the topology check. Every finite-dimensional normed vector
space over a complete valued field is complete and its coordinate maps
in a fixed basis are bounded. Induct on dimension: the span of all but
the last basis vector is complete by induction and hence closed; the
last vector has positive distance from it. Scaling this distance bounds
its coefficient, and the ultrametric inequality bounds the remaining
coefficients. This proves norm equivalence with a coefficient maximum
norm and completeness. Tensoring the bounded coordinate maps with \(A\)
and using the tensor norm upper bound identifies the tensor topology
with the finite-module topology on \(A\oplus Aa\), already complete.
This proves (4) for the actual unreduced algebra.

The involution on \(B\) is \(1\otimes\sigma\). The kernel of
\(\sigma-1\) on the \(K\)-vector space \(E\) is \(K\); tensoring this
kernel calculation over the field gives \(B^\sigma=A\). An invariant
unit and its inverse are both invariant, so invariant units descend as
units of \(A\). No field division by two is used. If an analytic element
has no zero on an affinoid spectrum, it is a unit: otherwise a maximal
ideal containing it and spectral surjectivity give a zero. This remains
valid with nilpotents.

The fibre over \(x\) is \(\mathcal M(\mathcal H(x)\otimes_K E)\).
Factoring the separable quadratic polynomial of \(a\) over
\(\mathcal H(x)\) gives a finite field or two field factors. Each finite
field factor has exactly one analytic point. Existence follows from the
spectrum of a finite affinoid field algebra. Uniqueness follows from
finite-dimensional norm equivalence: compare two extending absolute
values on \(z^m\) and take \(m\)th roots. Completeness follows from the
same norm argument. The involution is transitive on the factors: two
orbits would give an invariant nontrivial idempotent in
\((\mathcal H(x)\otimes E)^\sigma=\mathcal H(x)\). Thus a fibre is a
singleton fixed by \(\sigma\) or two interchanged points. This is not a
covering-space assertion.

The upstairs carrier is Hausdorff: points over different base points are
separated downstairs; points over one base point are separated inside
the Hausdorff scalar extension of an affinoid neighbourhood, with both
points in the inverse image of its ordinary interior. Inverse images of
compact sets are compact. Indeed cover the compact set by finitely many
interiors of affinoid neighbourhoods; its inverse image is closed in the
finite union of their compact scalar extensions.

We will use the following full-fibre shrinking fact. In a downstairs
affinoid neighbourhood \(V_0\) of \(x\), let \(N\subset V_{0,E}\) be a
relative ordinary open containing the entire fibre. The compact
complement has closed image in \(V_0\) not containing \(x\). Choose a
smaller strict affinoid neighbourhood \(V\) of \(x\) inside the
complement of this image and inside \(\operatorname{int}_X(V_0)\). Then
\(V_E\subset N\). Consequently finitely many strict norm inequalities or
unit conditions holding on a whole fibre can be retained simultaneously.
This applies to zero seminorm coefficients too; it does not replace the
algebra by its reduction.

\subsubsection{3.2. The restricted carrier and ambient
assumptions}\label{the-restricted-carrier-and-ambient-assumptions}

Let \(Z=X_E\) and let \(Z_0\) be its restricted \(K\)-analytic
structure. An \(E\)-affinoid algebra is \(K\)-affinoid under finite
restriction, since finite \(E/K\) is itself affinoid and the same Tate
presentations give the restriction. Rational domains in the common
Banach algebra use the same elements and inequalities. They give a
common affinoid neighbourhood basis and the same analytic structural
sheaf on ordinary opens. In particular unit germs and their compact
moment images coincide. The nullity criteria CLD 3.4.1 and 3.4.4 and the
ordinary moment descriptions identify the ordinary form sheaves, with
their supports, differentials and wedge products, on \(Z\) and \(Z_0\).
The calibrations differ by the factor two in Section 2(7). The
coefficient involution on (4) is a bounded \(K\)-algebra automorphism,
and the affinoid automorphisms glue to a genuine \(K\)-analytic
automorphism of \(Z_0\). This is the carrier on which it is used.

For \(X\) in the theorem, inverse images of ordinary affinoid interiors
give goodness upstairs. Strict presentations base change to strict
presentations. Separatedness follows by base change of the closed
diagonal. Boundarylessness follows from GR 2.11(5'). These assertions
also hold on ordinary opens. We make no claim that absolute analytic
interiors of arbitrary domains are preserved with equality.

We check purity at minimal components. In a strict affinoid
neighbourhood \(\mathcal M(A)\) of a pure \(n\)-dimensional space, every
minimal-prime component \(\mathcal M(A/P)\) has dimension \(n\). A point
of kernel \(P\) exists by spectral surjectivity, lies only on that
minimal component, and the local dimension formula proves this
assertion. Let \(J\) be the nilradical of \(A\), nilpotent by
noetherianity. The injection \[
 A/J\ \hookrightarrow\ \prod_P\operatorname{Frac}(A/P)
\] remains injective after tensoring with \(E\). Each field tensor is a
product of fields by separability. Hence \((A/J)\otimes E\) is reduced
and the nilradical of \(B\) is exactly \(J\otimes E\). The kernels of
the finitely many maps of this reduced algebra to these field factors
are primes with zero intersection. A minimal prime contains their
product, so contains one of them, and minimality makes it that kernel.
Its contraction to \(A\) is the associated minimal \(P\). The finite
integral injection \(A/P\hookrightarrow B/Q\) preserves dimension.

For completeness, the finite-domain dimension fact used in the last
sentence has the following algebraic proof. For a finite integral
injection of domains \(C\hookrightarrow D\) and a prime \(p\) of \(C\),
\(D_p\) is a nonzero finite \(C_p\)-module. If \(pD_p=D_p\), expressing
a finite generating list as \(p\)-linear combinations and taking the
determinant of \(I-M\) gives a unit annihilating \(D_p\), a
contradiction. Choose a maximal ideal containing \(pD_p\); its
contraction is \(pC_p\). This proves lying over. Quotient by a chosen
prime and repeat to prove going up. For incomparability, quotient by the
smaller of two primes with the same contraction and localize by the
nonzero elements of the base domain. The resulting finite integral
domain over a field is a field: a monic relation of least degree for a
nonzero element has a nonzero constant term and expresses its inverse. A
strictly larger prime could not remain after localization, and in a
domain a nonzero element cannot become zero under this localization.
Contraction of prime chains is therefore strict, and going up lifts
chains, proving equality of dimensions.

Every minimal component upstairs has dimension \(n\). The local
dimension formula proves purity upstairs. This reduction argument is
only a dimension check: all form measures still use the unreduced
analytic structure and its actual calibrated lengths. The checks repeat
at every finite quadratic tower stage.

\subsection{4. Scalar squares at the ramified
fold}\label{scalar-squares-at-the-ramified-fold}

\subsubsection{4.1. Normalization and every zero-coefficient
case}\label{normalization-and-every-zero-coefficient-case}

Work first near a singleton fibre \(\{y\}\) over \(x\) and take a finite
list of units \(f\) on \(V_E=\mathcal M(B)\). Express \(f=u+av\),
\(u,v\in A\). If \(v(x)=0\), then \(u(x)\ne0\). After a full-fibre
shrink, \(u\) is a downstairs unit and \(|av|,|\sigma(a)v|<|u|\) on all
of \(V_E\). Consequently \[
 t_f=t_{\sigma f}=t_u,\qquad \delta_f:=t_f-t_{\sigma f}=0     \tag{5}
\] as ordinary functions and, on differentiating their equality, as
ordinary forms. Nilpotent \(v\) is included in this case.

If \(v(x)\ne0\), shrink so \(v\) is a downstairs unit and put \[
 b=u/v,\qquad F=(a+b)/c=f/(cv).
\] Then \(F,\sigma F\) are units and \[
 F+\sigma F=1,\qquad
 t_F-t_{\sigma F}=\delta_f.                                \tag{6}
\] Indeed \(\sigma F=-(\sigma(a)+b)/c\) and \(|\sigma c|=|c|\). For
\(U=t_F\), \(V=t_{\sigma F}\) the ultrametric inequality puts the moment
image in the standard tripod \[
 \{U=V\leq0\}\ \cup\ \{U\geq0,V=0\}\ \cup\ \{V\geq0,U=0\}.
                                                                    \tag{7}
\] Thus, with the downstairs unit \(N_F=F\sigma F\) and \(S_f=t_{N_F}\),
\[
 \delta_f^2=q(S_f\circ\pi),\qquad q(s)=(\max(s,0))^2.        \tag{8}
\] This is an identity at every analytic point, not just on one
skeleton. The three rays prove it separately. The function \(q\) is
\(C^1\) but not \(C^2\) at zero, so (8) is not asserted to be ordinary
exact scalar descent. At the singleton fibre \(|F|=|\sigma F|\geq1\). If
the inequality is strict, retain both magnitudes \(>1\) by shrinking;
then \(\delta_f=0\) throughout. If equality holds, (8) handles the
vertex without dividing by \(\delta_f\).

Perform these constructions for the original units \(f_i\) and all their
pair products \(f_if_j\) for \(i<j\). A single shrink suffices because
the list is finite. In the zero case assign \(R_{f,\varepsilon}=0\)
below; otherwise retain its unit \(N_F\) and coordinate \(S_f\).

\subsubsection{4.2. A smooth fold with controlled differentiated
error}\label{a-smooth-fold-with-controlled-differentiated-error}

Fix an even, nonnegative \(C^\infty\) function \(\rho\) supported in
\([-1,1]\) with integral one; for example normalize \(\exp(-1/(1-u^2))\)
on \(|u|<1\), extended by zero. Define \[
 q_\varepsilon(s)=\int\rho(u)q(s-\varepsilon u)\,du
 \quad(0<\varepsilon\leq1).                               \tag{9}
\] This is smooth. Since \(q'(s)=2\max(s,0)\) is Lipschitz with constant
two and \(\int u\rho(u)du=0\), the integral remainder formula gives \[
 |q_\varepsilon-q|\leq\varepsilon^2\int u^2\rho(u)du,
 \quad |q'_\varepsilon-q'|\leq2\varepsilon\int|u|\rho(u)du.
                                                                    \tag{10}
\] The weak derivative of \(q'\) is \(2\mathbf1_{s>0}\), with no atom at
zero because \(q'\) is continuous. Directly differentiating its
convolution gives \[
 0\leq q''_\varepsilon\leq2,\qquad
 q''_\varepsilon(s)=2\mathbf1_{s>0}\quad\text{if }|s|>\varepsilon.
                                                                    \tag{11}
\] These assertions also follow by splitting the integral in (9) at
\(u=s/\varepsilon\). No analytic current has been differentiated to
obtain them.

Set \(R_{f,\varepsilon}=q_\varepsilon(S_f)\) in the nonzero case. These
are ordinary smooth downstairs functions. Define \[
 Q_{ii,\varepsilon}=R_{f_i,\varepsilon},\qquad
 Q_{ij,\varepsilon}
  =\tfrac12(R_{f_if_j,\varepsilon}
                   -R_{f_i,\varepsilon}-R_{f_j,\varepsilon}),\quad i<j,
                                                                    \tag{12}
\] and \(Q_{ji,\varepsilon}=Q_{ij,\varepsilon}\). The target upstairs is
\(Q_{ij}=\delta_i\delta_j\) because
\(\delta_{f_if_j}=\delta_i+\delta_j\).

Here is the exact topology in which their derivatives converge. Add to
one fixed upstairs moment tuple all the original units, conjugates,
normalized \(F\) and \(\sigma F\), and the pulled-back downstairs norm
units. On any compact domain its moment image is a finite polyhedral
set. Refine by the hyperplanes \(S_f=0\). On the relative interior of
each cell and for each \(f\), either \(S_f>0\), \(S_f<0\), or \(S_f\) is
identically zero. In the positive case \(\delta_f=\pm S_f\) on that
cell; in the negative case \(\delta_f=0\); in the zero case both
coordinates in (7) and \(\delta_f\) are identically zero. These are
affine identities on the cell, so can be differentiated along its
tangent space.

It follows from (10)--(11) that the restrictions to cell tangents of \[
 R_{f,\varepsilon}\circ\pi,\quad d'(R_{f,\varepsilon}\circ\pi),
 \quad d''(R_{f,\varepsilon}\circ\pi),\quad
 d'd''(R_{f,\varepsilon}\circ\pi)                          \tag{13}
\] converge to those of \(\delta_f^2\) and its corresponding exterior
derivatives at almost every point of every top cell. Their coefficients
in orthonormal tangent bases stay bounded on the fixed compact image.
Zeroth and first errors are \(O(\varepsilon^2)\) and \(O(\varepsilon)\).
The second error is uniformly bounded and is zero outside
\(|S_f|\leq\varepsilon\) on a cell where \(S_f\) is nonconstant. On a
cell where \(S_f\) is constant, all its tangential derivatives are zero;
this includes an entire cell contained in the fold. Thus such a cell
does not contribute a second derivative error, even though
\(q''_\varepsilon(0)\) need not converge to either one-sided value. The
same assertions hold for (12) by polarization.

These are tangential convergence statements on actual moment cells, not
ambient \(C^2\) convergence at the fold.

\subsection{5. Actual calibration and the thin-strip
bound}\label{actual-calibration-and-the-thin-strip-bound}

We record the precise measure argument that will turn the topology just
described into differentiated compact pairings. Let \(g\) be a fixed
unit moment tuple of length \(r\) on a compact domain in a pure
\(n\)-dimensional analytic space, and put \[
 \xi=\sum_{k=1}^r d't_k\wedge d''t_k,\qquad
 \zeta=\xi^n/n!
\] with the standard positive top-degree ordering. If \(r<n\), every top
form in that tuple vanishes. Otherwise, for a tangent \(n\)-volume
represented by an \(r\times n\) matrix \(V\), write \(v_I=\det(V_I)\),
\(|I|=n\). A coefficient top form \(\eta=\sum H_{IJ}d't_I\wedge d''t_J\)
contracts, up to the common ordering convention, to
\(\sum H_{IJ}v_Iv_J\) times its positive canonical calibration. The form
\(\zeta\) contracts to \(\sum_Iv_I^2\) times the same calibration. Thus
if \(|H_{IJ}|\leq M\), \[
 |\eta_{\mathrm{Rad}}|
      \leq {r\choose n}M\,\zeta_{\mathrm{Rad}}.             \tag{14}
\] Indeed \((\sum_I|v_I|)^2\leq {r\choose n}\sum_Iv_I^2\). The same
bound applies using bounded tangential coefficients instead: choose
orthonormal frames on the finitely many cell tangent spaces; then
\(\sum_Iv_I^2=1\) by the determinant expansion of \(\det(V^tV)\), and
their coefficient contraction gives a fixed finite bound against
\(\zeta\). This is a cellwise bound on the actual supporting skeleton,
so positivity gives a total-variation bound on the carrier, not merely a
bound on a signed pushforward with possible fibre cancellation.

The measure \(\zeta_{\mathrm{Rad}}\) is finite on the compact domain by
the calibrated contraction formula and CLD 5.10.1(c,d). Its moment
pushforward on each top cell has a constant finite positive density
relative to the Euclidean volume chosen there. Actual generic lengths
are in this density; the same density multiplies both sides of (14). No
reduction, fibre point count, or probability normalization is used.

For a nonconstant affine \(S\) on one compact top polytope \(C\), the
\(\zeta\)-measure of \(\{|S|\leq\varepsilon\}\cap C\) is
\(O(\varepsilon)\). To see this directly, use an orthonormal coordinate
on its affine span parallel to the nonzero gradient of \(S\). The strip
has width \(2\varepsilon/\|\nabla_C S\|\) in that coordinate. All its
orthogonal cross sections lie in one fixed box containing \(C\), so
Fubini bounds its volume by that width times the volume of the box's
cross section. Multiply by the fixed calibration density and sum over
the finitely many cells. Cells where \(S\) is constant require no strip
estimate for (13), since \(dS\) restricts to zero. Lower dimensional
cell faces have zero top measure. Hyperplane intersections in a
nonconstant cell have zero top measure as well.

Consequently, for any finite collection of fold coordinates and any top
forms whose tangential coefficient error is bounded by \[
 C\varepsilon+C\sum_{f,\,S_f|_C\text{ nonconstant}}
       \mathbf1_{\{|S_f|\leq\varepsilon\}},                 \tag{15}
\] their Radon total-variation error on the fixed compact domain tends
to zero, in fact is \(O(\varepsilon)\). This includes cells entirely in
the fixed locus. Any conductor stratum contained in a lower dimensional
cell has zero top measure; a top cell contained in a fold was handled by
its zero tangential derivative. Any further top cells, singular or
otherwise, are included with their actual calibration. We do not assume
smoothness of the analytic space, or that a singular locus can simply be
discarded. This exhausts the calibrated contributions, including
nonreduced multiplicities.

\subsection{6. Orbital one-forms and the cancellation of high
derivatives}\label{orbital-one-forms-and-the-cancellation-of-high-derivatives}

\subsubsection{6.1. The tripod identity in the ordinary
sheaf}\label{the-tripod-identity-in-the-ordinary-sheaf}

Let \(F,G,e\) be analytic units with \(G=F+e\). Put \(x=t_F\), \(z=t_G\),
\(r=t_e\), \(U=x-r\), \(V=z-r\). The relation \(G/e=1+F/e\) puts
\((U,V)\) in the three-ray set (7). For either \(d=d'\) or \(d=d''\),
the polynomial one-form \(U\,dV-V\,dU\) restricts to zero on each ray:
on the diagonal the terms cancel; on an axis the zero coordinate and its
tangential differential vanish. At the vertex both coefficients are
zero. Its restriction to every compact moment image is therefore zero.
CLD 3.4.4, using this actual unit moment on an ordinary neighbourhood,
proves its nullity in the ordinary form sheaf. Expanding gives \[
 x\,dz-z\,dx=r(dz-dx)+(x-z)\,dr.                           \tag{16}
\] This proof is not differentiation of pointwise values at isolated
points, and is not an abstract tensor model.

For the original indices with \(v_i(x)\ne0\) retain \(F_i=(a+b_i)/c\).
If \(e_{ij}=b_j-b_i\) has zero value at \(x\), shrink so
\(|e_{ij}/c|<|F_i|,|\sigma F_i|\) on the full upstairs neighbourhood.
Then the corresponding conjugate logarithms for \(i\) and \(j\) agree,
and \(\delta_i=\delta_j\), including their differentials. The orbital
one-form is zero. If \(e_{ij}(x)\ne0\), shrink so \(e_{ij}\) is a
downstairs unit and apply (16) to \(F_j=F_i+e_{ij}/c\) and to its
conjugate. The two difference units \(e_{ij}/c\) and \(-e_{ij}/c\) have
the same logarithm \[
 r_{ij}=t_{e_{ij}}+\log|c|,
\] an ordinary downstairs unit logarithm plus a real constant. Let
\(S_i=t_{F_i\sigma F_i}\), \(S_j=t_{F_j\sigma F_j}\) downstairs. Adding
(16) for the conjugate pair gives the exact ordinary identity \[
 \pi^*O_{ij}^{d}=\delta_i\,d\delta_j-\delta_j\,d\delta_i,
                                                                    \tag{17}
\] where \[
 O_{ij}^{d}=2r_{ij}(dS_j-dS_i)+2(S_i-S_j)\,dr_{ij}
                         -S_i\,dS_j+S_j\,dS_i.             \tag{18}
\] To check the factor, the sum of the left sides of (16) is
\(\tfrac12(S_i dS_j-S_j dS_i+
\delta_i d\delta_j-\delta_j d\delta_i)\). If one original \(v_i\)
vanishes at \(x\), (5) makes the orbital expression zero and we set
\(O_{ij}^d=0\). In the zero-\(e_{ij}\) case also set it zero. Set
\(O_{ii}^d=0\), \(O_{ji}^d=-O_{ij}^d\). All choices are finite and hold
on one smaller neighbourhood. No vanishing coefficient or logarithmic
difference is divided out.

The coefficients of (18) are affine in the fixed unit-log coordinates
\(S_i,S_j,r_{ij}\). Hence \[
 d'd''O_{ij}^{d}=0,\qquad d'O_{ij}^{d},d''O_{ij}^{d}
        \text{ have constant coefficient representations}.          \tag{19}
\] This statement is about superforms in those moment coordinates.

\subsubsection{6.2. Smoothed generators and their full differentiated
families}\label{smoothed-generators-and-their-full-differentiated-families}

Define ordinary downstairs one-forms \[
 L_{ij,\varepsilon}^{d}
      =\tfrac12(dQ_{ij,\varepsilon}+O_{ij}^{d}).             \tag{20}
\] Their upstairs targets are \(L_{ij}^{d}=\delta_i\,d\delta_j\). By
(17) and the product rule for \(d(\delta_i\delta_j)\), the tangential
convergence in (13) proves convergence of (20),
\(d'L_{ij,\varepsilon}^{d}\), and \(d''L_{ij,\varepsilon}^{d}\) to their
targets, with bounded tangent coefficients on each fixed compact chart.
Moreover \[
 d'd''L_{ij,\varepsilon}^{d}=0                              \tag{21}
\] \textbf{exactly}: the contribution of \(dQ_{ij,\varepsilon}\)
vanishes by \((d')^2=(d'')^2=0\) and \(d'd''=-d''d'\), and that of \(O\)
by (19). The upstairs target has the same identity since unit logarithms
have zero \(d'd''\) as forms.

Use the two-form generators \[
 P_{ij,\varepsilon}^{++}=d'L_{ij,\varepsilon}^{d'},\quad
 P_{ij,\varepsilon}^{--}=d''L_{ij,\varepsilon}^{d''},\quad
 P_{ij,\varepsilon}^{+-}=d'L_{ij,\varepsilon}^{d''}.          \tag{22}
\] Their respective targets are \(d'\delta_i\wedge d'\delta_j\),
\(d''\delta_i\wedge d''\delta_j\), and \(d'\delta_i\wedge d''\delta_j\),
with ordered indices in the last case. They converge tangentially with
bounded coefficients by (13). Each two-form in (22) is closed for
\textbf{both} \(d'\) and \(d''\): one closure follows from the
square-zero identity and the other from (21). Its target is also closed
for both.

Thus the family consisting of a generator and its \(d',d'',d'd''\)
derivatives converges tangentially, with uniformly bounded coefficients,
for each \(Q,L,P\). Errors of its entries are bounded by a constant
times the form of (15); for \(Q\) and \(L\) themselves the error is
\(O(\varepsilon)\) without a strip term. This follows from (10)--(13)
and the exact cancellations (19)--(22). There is no use of third or
fourth derivatives of \(q_\varepsilon\). Such derivatives are unbounded
near the fold and would invalidate a naive coefficient \(C^2\) argument
for (22).

\subsubsection{6.3. Every invariant polynomial
superform}\label{every-invariant-polynomial-superform}

Put \(s_i=t_{f_i}+t_{\sigma f_i}\) and
\(\delta_i=t_{f_i}-t_{\sigma f_i}\). The \(s_i\) descend exactly as norm
logarithms, as do their prime differentials. In the free real
polynomial/exterior algebra in \(s_i,\delta_i\) and their prime
differentials, simultaneous conjugation fixes the \(s\) symbols and
changes the sign of every \(\delta\) symbol, counting scalars and
differentials alike. An invariant monomial contains an even number of
anti-invariant symbols. This follows by comparing coefficients in the
free monomial basis; it does not impose symmetry on mixed forms. Pair
these symbols, moving differentials with their graded signs. Every pair
is one of \[
 \delta_i\delta_j,\quad \delta_i d'\delta_j,\quad
 \delta_i d''\delta_j,\quad
 d'\delta_i\wedge d'\delta_j,\quad
 d''\delta_i\wedge d''\delta_j,\quad
 d'\delta_i\wedge d''\delta_j,
\] or the negative of an ordered pair. These are precisely the targets
of (12), (20), (22). Replace the pairs by their smoothed downstairs
generators and leave the invariant norm symbols unchanged. This gives an
ordinary downstairs form \(A_{P,\varepsilon}\) for every invariant
polynomial superform \(P\) in \((t_f,t_{\sigma f})\).

For a fixed polynomial, any such pairing decomposition suffices; we make
no independence-of-decomposition assertion at finite \(\varepsilon\).
All unit lists, shrinkings, and moments are independent of its degree.
The Leibniz rules show that the differentiated family of a product is a
finite sum of products of entries in its factors' families. Equations
(21)--(22) therefore ensure that its tangent coefficients are bounded
and converge, and its top differentiated error satisfies (15), with a
constant allowed to depend on \(P\). This proves, on every fixed compact
moment domain and also after multiplication by any fixed ordinary smooth
cutoff \(\chi\), \[
 \big\|(d'd''(\chi(P-\pi^*A_{P,\varepsilon})))_{\mathrm{Rad},E}
       \big\|_{\mathrm{TV}}\longrightarrow0.                \tag{23}
\] Here the expression is a top form when \(P\) has bidegree
\((n-1,n-1)\). The proof uses the \(d',d'',d'd''\) families just listed,
the bounded corresponding family of \(\chi\), and Section 5. It treats
alternating pure-prime and ordered mixed terms, not only scalar squares
or symmetric Hessians. For \(n=1\) it simply specializes to scalar
tests.

\subsection{7. All compact ordinary tests with fixed
support}\label{all-compact-ordinary-tests-with-fixed-support}

Let \(\beta\in A_c^{n-1,n-1}(X_E)\) be invariant under \(\sigma\). We
prove that there are \(\alpha_j\in A_c^{n-1,n-1}(X)\) and one compact
\(D\subset X\) with \(\operatorname{supp}\alpha_j\subset D\) and \[
 \big\|(d'd''(\beta-\pi^*\alpha_j))_{\mathrm{Rad},E}
       \big\|_{\mathrm{TV}}\to0,
 \qquad \operatorname{supp}\subseteq\pi^{-1}(D).            \tag{24}
\] This asserts only the differentiated approximation required for the
current equation, not exact ordinary descent or ambient \(C^2\) density
of the ramified descendants.

Near a singleton fibre, choose an ordinary moment representation of
\(\beta=f_{\mathrm{trop}}^*\gamma\) on a neighbourhood of its point,
intersect that neighbourhood with its conjugate, and shrink downstairs
until \(V_E\) lies in it. The Euclidean coefficient extension in Section
2(3), multiplied by a Euclidean cutoff equal to one near the compact
image, lets us take \(\gamma\) smooth on the whole coordinate space. The
coefficient superform \[
 G=\tfrac12(\operatorname{pr}_1^*\gamma+
                           \operatorname{pr}_2^*\gamma)               \tag{25}
\] on the doubled tuple \((f,\sigma f)\) is swap-invariant and pulls
back to \(\beta\). Its invariance includes its differential basis
factors. Carry out Sections 4--6 for this fixed tuple and its pair
products. Choose a swap-stable box containing the compact moment image
in its interior. There are invariant polynomial superforms \(P_j\) whose
coefficient errors from \(G\) converge in \(C^2\) on that box.

Here is a direct proof of the polynomial approximation statement. Extend
each smooth coefficient near the box by a compactly supported smooth
coefficient agreeing there. Convolve with a Gaussian. Its derivatives
through order two converge uniformly on the box: use uniform continuity
of the compactly supported derivatives, split the integral into a small
ball and its tail, and differentiate under the integral. For fixed width
the Gaussian kernel and its derivatives through order two are uniform
limits of their Taylor polynomials on the compact difference of the box
and coefficient support. Integrate these finite polynomials against the
coefficient to obtain polynomial approximants with the same derivative
control. Choose decreasing widths and sufficient truncation degrees. Do
this for the finitely many coefficient functions, then average the
superform with its coordinate-swapped pullback. This only permutes
differential indices with fixed signs and preserves \(C^2\) control on
the swap-stable box. It proves the assertion without an imported
approximation theorem.

Near a two-point fibre choose disjoint ordinary neighbourhoods \(N_1\)
and \(N_2=\sigma N_1\), with a moment representation of \(\beta\) on
\(N_1\). Shrink downstairs until \(V_E\subset N_1\cup N_2\). Glue each
unit as \(f_i\) on the first piece and \(\sigma f_i\) on the second. The
pieces are disjoint opens covering the affinoid; the structural sheaf
glues these units and their inverses. They are invariant and therefore
are units downstairs by (4). Invariance of \(\beta\), including the
differential factors, gives its exact descendant in these units. Use
that same descendant at every approximation stage. This proof does not
assume a local topological covering.

The compact image \(C=\pi(\operatorname{supp}\beta)\) is covered by
finitely many ordinary interiors \(U_i\) of these downstairs
neighbourhoods. Take fixed cutoffs \(\chi_i\) from the compact finite-cover
construction of Section 2, using CLD 15.2.5, compactly supported in \(U_i\), with sum one near \(C\), and set
\(D=\bigcup_i\operatorname{supp}\chi_i\). This is compact and contains
\(C\). On singleton charts let
\(b_{i,j,\varepsilon}=A_{P_{i,j},\varepsilon}\); on two-point charts let
it be the exact descendant. Set \[
 \alpha_{j,\varepsilon}=
        \sum_i\widetilde{\chi_i b_{i,j,\varepsilon}},        \tag{26}
\] where tilde means extension by zero from \(U_i\). These are ordinary
smooth compact forms. Indeed outside the closed compact support of
\(\chi_i\), including at every point outside \(U_i\), an ordinary open
neighbourhood makes the product zero, so it glues smoothly to zero. No
boundary distributions arise from this extension.

The global error equals the sum of these fixed cutoff local errors: \[
 \beta-\pi^*\alpha_{j,\varepsilon}
     =\sum_i(\chi_i\circ\pi)
                  (\beta-\pi^*b_{i,j,\varepsilon}).         \tag{27}
\] The sum-one identity is required only near the support of \(\beta\);
off that support all terms involving \(\beta\) are zero. Cover
\(\pi^{-1}(\operatorname{supp}\chi_i)\) by finitely many compact
upstairs moment neighbourhoods lying within its local construction and
providing fixed moment descriptions of \(\chi_i\circ\pi\). Combine these
moments with all the moments used in Sections 4--6. Every chart, compact
domain and cutoff is now fixed, independent of \(j\) and
\(\varepsilon\).

Split each singleton error into \(\beta-P_{i,j}\) and
\(P_{i,j}-\pi^*b_{i,j,\varepsilon}\). For the first, coefficient \(C^2\)
convergence in fixed original upstairs moment coordinates, the product
rule with the fixed cutoff, and (14) give differentiated total-variation
convergence to zero as \(j\to\infty\). This use of \(C^2\) is
legitimate; it concerns Euclidean polynomial approximation to (25), not
the fold. For the second, (23) gives convergence to zero as
\(\varepsilon\to0\) for each \textbf{fixed} \(j\). The finite number of
terms and charts permits a choice \(\varepsilon_j\in(0,1/j)\) for which
their summed differentiated total variation is below \(1/j\). No bound
uniform in polynomial degree is needed. Choose
\(\alpha_j=\alpha_{j,\varepsilon_j}\). The triangle inequality proves
(24) with support in the fixed compact \(\pi^{-1}(D)\).

For every continuous finite \(h\), boundedness of \(h\circ\pi\) on this
compact set therefore implies \[
 B_{h\circ\pi}(\beta-\pi^*\alpha_j)\longrightarrow0.        \tag{28}
\] Thus the relaxed topology explicitly controls differentiated compact
pairings for every such multiplier, including all fold and calibrated
length contributions. Scalar \(C^0\) density alone would not establish
(28).

\subsection{8. Averaging, the two field categories, and the
theorem}\label{averaging-the-two-field-categories-and-the-theorem}

On the restricted carrier \(Z_0\) in Section 3.2 the involution is a
\(K\)-analytic isomorphism of calibrated degree one. CLD 15.4.3 gives,
for a compact smooth top form \(\omega\), \[
 \sigma_*((\sigma^*\omega)_{\mathrm{Rad},K})
                      =\omega_{\mathrm{Rad},K}.
\] The same-carrier calibration relation from CLD 14.4.3 and equality of
ordinary sheaves gives
\(\omega_{\mathrm{Rad},K}=2\omega_{\mathrm{Rad},E}\). Cancelling two
proves \[
 \sigma_*((\sigma^*\omega)_{\mathrm{Rad},E})
                      =\omega_{\mathrm{Rad},E}.             \tag{29}
\] Since \(h\circ\pi\) is invariant and the involution commutes with
ordinary exterior differentials, (29) shows
\(B_{h\circ\pi}(\sigma^*\beta)=B_{h\circ\pi}(\beta)\) for every compact
upstairs test. It may be replaced by its real average
\(\beta^{\mathrm{av}}=(\beta+\sigma^*\beta)/2\), with the same pairing.
Its support lies in a compact finite saturation. This averaging uses
neither field division by two nor an unproved \(E\)-analytic involution.

For a downstairs compact test \(\alpha\), its pullback is compact by
Section 3.1. Smooth functoriality and the \textbf{scalar-projection}
identity (3), which has no factor two, give \[
 \begin{split}
 B_{h\circ\pi}(\pi^*\alpha)
 &=\int_{X_E}(h\circ\pi)(\pi^*d'd''\alpha)_{\mathrm{Rad},E}\\
 &=\int_X h(d'd''\alpha)_{\mathrm{Rad},K}=B_h(\alpha).
 \end{split}                                                \tag{30}
\] These are finite signed measures supported on a fixed compact set;
integration against a bounded continuous multiplier is legitimate. No
continuous-current pullback is constructed by notation.

Suppose \(d'd''[h]=0\). Apply (24)--(28) to \(\beta^{\mathrm{av}}\).
Every approximating pairing is zero by (30), so
\(B_{h\circ\pi}(\beta^{\mathrm{av}})=0\) and hence the pairing with
every upstairs test is zero. This proves the new forward implication.
Conversely, if the upstairs equation vanishes, (30) for every downstairs
test proves the reverse implication. This proves (1). Ordinary opens
retain the ambient assumptions. The checks in Section 3.2 justify
iteration through every finite separable quadratic tower; the composite
need not be Galois. No passage to an infinite completion has been made.

\subsection{9. The genuinely ramified example and comparison with
P06.2}\label{the-genuinely-ramified-example-and-comparison-with-p06.2}

The polynomial \(T^2-2\) is irreducible over \(\mathbb Q_2\): a square
in \(\mathbb Q_2^\times\) has even integral \(2\)-adic valuation,
whereas \(v_2(2)=1\). It is separable, and its other root is
\(-\sqrt2\). For \(a=\sqrt2\), \(c=2\sqrt2\), \[
 |a|=2^{-1/2},\qquad |c|=2^{-3/2}.
\] No unit in this extension satisfies the earlier residue-separating
condition. Indeed write \(z=x+y\sqrt2\), \(x,y\in\mathbb Q_2\). The
nonzero summands have valuations in \(\mathbb Z\) and \(\mathbb Z+1/2\),
respectively, so cannot cancel and \(v_2(z)=\min(v_2(x),v_2(y)+1/2)\).
If \(z\) is a unit this forces \(v_2(x)=0\) and \(v_2(y)\geq0\) (or
\(y=0\)). Thus \(|z-\sigma z|=|2y\sqrt2|\leq2^{-3/2}<1\) (or is zero).
The field falls strictly outside the parent condition
\(|z|=|z-\sigma z|=1\). Equations (6)--(11), not the old norm-one trace
formula, are what treat it here. The argument works unchanged for every
separable quadratic extension of every complete nontrivially valued
field.

Relative to the catalogue, this is automatic transport for arbitrary
pure-dimensional strict analytic spaces under these extensions,
including nonproduct singular and nonreduced spaces. It does not assume
an upstairs zero-current equation, a toric invariant presentation, a
product chart, a proper full-rank moment factor, or ordinary exact
invariant form descent. It is a new forward gate beyond the parent
restricted field condition; descent alone was already an earlier scope.

It is \textbf{not} a complete solution of original P06.2. We have not
proved:

\begin{itemize}
\tightlist
\item
  recognition of every continuous zero-current function as classically
  pluriharmonic on arbitrary good separated strict boundaryless spaces;
\item
  ordinary two-sign positive-current regularization, either on arbitrary
  germs or as a consequence of (1);
\item
  scalar-current ascent for general finite separable extensions outside
  finite quadratic towers, infinite algebraic completions, or
  transcendental analytic extensions;
\item
  the universal classical curve-probe bridge over all analytic
  extensions;
\item
  exact ordinary smooth descent of every invariant form in the ramified
  case, ordinary smoothness of \(h\), or a harmonic-approximant theorem.
\end{itemize}

The theorem transports a \textbf{zero} equation; its approximation is
not asserted to preserve positivity of tests, and hence is not
automatically a theorem transporting arbitrary positive-current
inequalities. Ordinary CLD regularization, weak smoothness, and local
real-PL recognition remain distinct. No analytic counterexample is
claimed. The catalogue's non-strict arbitrary-radius boundary cutoff is
outside the original strict ambient scope and is not used. P05.2 is not
addressed.

\subsection{10. Evidence and
reproducibility}\label{evidence-and-reproducibility}

This file is the authoritative full argument; all new supporting lemmas
used by it are proved here. The complete standalone typeset companion is
\texttt{P06.2-partial.tex} and \texttt{P06.2-partial.pdf}.
\texttt{self-audit.md} checks the dependency gates and limitations.
\texttt{research-notes.md} preserves the failed routes and concrete
remaining obligations. \texttt{status.json} has exactly the five
requested keys. \texttt{input.json} retains the supplied task and
assignment.

Run \texttt{python3\ build-and-check.py} from this workspace. It
performs only finite exact diagnostic identities and two installed
pdfLaTeX builds, records all commands and outputs, and writes
source/build/artifact hash evidence. Source snapshots and fallible
parent copies are in \texttt{sources/}; their origins and hashes are in
\texttt{source-manifest.json}. \texttt{evidence/primary-passages.txt}
retains the inspected source passages, and
\texttt{evidence/environment.json} records the environment. The
diagnostic program checks cancellations algebraically and fold values
exactly in a finite smoothing surrogate; it does not verify the analytic
theorem. Typesetting success is not mathematical verification. No
installation, additional model/agent, external message, old-artifact
write, or heavy computation was used. The exact partial theorem was independently reviewed through Xean native
Codex; the full PASS report and all source bindings are preserved.
No formal verification or priority certification is claimed. The saved
publication source was opened in the native desktop editor and its
built-in compiler returned success; two local pdfLaTeX passes exported
the same source as a PDF. Compilation checks typesetting only.

\subsection{Bibliography}\label{bibliography}

\begin{itemize}
\tightlist
\item
  A. Chambert-Loir and A. Ducros, \emph{Formes différentielles réelles
  et courants sur les espaces de Berkovich}, arXiv:1204.6277v2, 25 July
  2025. Exact inputs and hypotheses are stated in Section 2.
\item
  W. Gubler and J. Rabinoff, \emph{Classically psh and pluriharmonic
  functions on Berkovich spaces}, arXiv:2506.13548v2, 17 September 2025,
  Section 2.11(5').
\item
  Supplied original P06.2 dossier and preceding restricted-quadratic
  manuscript. Fallible mathematical comparison data only; no imported
  theorem assumptions or verification receipts.
\end{itemize}

\end{document}
